\documentclass{article}
\usepackage{graphicx} % Required for inserting images
\usepackage[a4paper, total={16cm, 23cm}]{geometry}
\usepackage{parskip}
\usepackage{amsmath} % Equations

\title{Beispieltest}
\author{Platz für Hinweise}
\date{Name, Punkte o. Ä.}

\newcommand{\exercise}[2]{\item \expandafter\ifx\expandafter\relax
	\detokenize{#1}\relax\else #1 \\ \fi #2 \\}

\newcommand{\inputline}{\rule{30pt}{0.4pt}}

\begin{document}

\maketitle

\vspace{5pt}
\hrule
\vspace{10pt}

\begin{enumerate}

    \exercise{}{
    If there are five apples and you take away three, how many do you have?

\inputline
}\exercise{}{
    One grandmother, two mothers, two daughters and one granddaughter go to
the cinema and buy one ticket each. How many tickets do they have to buy
in total?

\inputline
}\exercise{}{
    Einstein's most famous formula, relating Energy \(E\) and mass \(m\) via
the speed of light \(c\), reads \(E=\)\inputline.
}\exercise{}{
    Give a factor of 42: \inputline
}\exercise{}{
    36 = \inputline \(\times\) \inputline
}\exercise{}{
    What is the answer to the Ultimate Question of Life, The Universe, and
Everything?

\inputline
}\exercise{}{
    Free lunch!
}\exercise{}{
    4 divided by 2 is \inputline with remainder \inputline.
}\exercise{}{
    9 \(\times\) 10 = \inputline
}\exercise{}{
    The Pythagorean Theorem reads \inputline.
}\exercise{}{
    A half is a third of it. What is it?

\inputline / \inputline
}\exercise{}{
    A `Sextillion' equals \inputline.
}\exercise{
    You know that \(2 + 2 = 2 \times 2\).
}{
    Find a set of three different integers whose sum is equal to their
product.

\inputline
}\exercise{}{
    The square root of 100 is \inputline.
}

<<exercises>>

\end{enumerate}

\end{document}